% Generated draft: requires mathematical and source review.
\documentclass[11pt]{article}
\usepackage[margin=25mm]{geometry}
\usepackage{fontspec}
\setmainfont{Latin Modern Roman}
\usepackage{amsmath,amssymb,mathtools,unicode-math}
\setmathfont{Latin Modern Math}
\usepackage{enumitem,needspace,xurl,hyperref}
\setlength{\parindent}{0pt}
\setlength{\parskip}{5pt}
\setlength{\emergencystretch}{3em}
\newcommand{\N}{\mathbb{N}}
\newcommand{\Z}{\mathbb{Z}}
\newcommand{\Q}{\mathbb{Q}}
\newcommand{\R}{\mathbb{R}}
\newcommand{\C}{\mathbb{C}}
\newcommand{\F}{\mathbb{F}}
\newcommand{\A}{\mathbb{A}}
\newcommand{\PP}{\mathbb{P}}
\newcommand{\E}{\mathbb{E}}
\newcommand{\eps}{\varepsilon}
\newcommand{\dd}{\,\mathrm{d}}
\begin{document}
\section*{UnsolvedMath ID: 11000030}
\textbf{Source problem code:} catalogue entry
\subsection{Definitions and mathematical statement}
Let \(H\le\operatorname{Mod}_g\) be a finitely presented subgroup of the source's mapping class group. The conjugacy problem in \(H\) is to decide, from words representing \(x,y\in H\) in a fixed finite generating set, whether
\[
 \exists h\in H:\quad hxh^{-1}=y.
\]
The conjugator must lie in \(H\), not merely in \(\operatorname{Mod}_g\). A solvable problem means an algorithm terminating on every such pair.

Question 3.9 asks whether every finitely presented \(H\) has solvable conjugacy problem, and separately whether every such \(H\) is combable and whether it is automatic. In a usual combing convention, choose a word path from the identity to each element, with a uniform fellow-traveler bound for paths to adjacent endpoints. Automaticity additionally requires a regular language of representatives and the appropriate finite-state multiplier/equality relations. The source does not specify synchronous/asynchronous or other combing conventions, so those definitions need confirmation; no biautomaticity requirement is added.

The questions are about all finitely presented subgroups, not arbitrary finitely generated subgroups. They do not specify a single uniform algorithm taking a subgroup presentation as an additional input. The genus and surface conventions are omitted.

The context notes that solvable word problem passes to finitely generated subgroups, but solvable conjugacy problem need not: it cites subgroups of \(\mathrm{GL}(n,\Z)\) with unsolvable conjugacy problem, reference [Mi]. It also states that \(\operatorname{Mod}_g\), like \(\mathrm{GL}(n,\Z)\), has only finitely many conjugacy classes of finite subgroups. These observations do not answer the intrinsic-subgroup questions. The supplied summary's assertion that mapping class groups are not biautomatic for \(g\ge2\) has no supplied reference or verification and is not made a hypothesis or a conclusion of this draft.
\subsection{Short English statement}
For every finitely presented subgroup of a mapping class group, ask separately whether its intrinsic conjugacy problem is decidable, whether it is combable and whether it is automatic. The source gives comparison examples and missing references, but no answer or uniform subgroup-input algorithm.
\subsection{Sources}
\begin{itemize}
\item Local UnsolvedMath registry, permanent ID \(11000030\), exact source code \texttt{catalogue entry}, Question 3.9, PDF page 34. No individual proposer or proposed year is supplied.
\item Farb, edited volume \emph{Problems on Mapping Class Groups and Related Topics}: \url{https://www.math.uchicago.edu/~farb/papers/mcgbook.pdf}. The volume contains 329 explicitly typeset numbered items across contributed chapters, with repeated numbering distinguished by PDF page.
\item Internal reference [Mi] for linear-group subgroups with unsolvable conjugacy problem; the bibliography entry is absent and no author or title is guessed.
\item Extraction pdf-font-structure, accessed 2026-07-29. Source evidence NEEDS\_REVIEW presents an open problem and requires a current primary-literature audit. Supplied OPEN-TRIAGE summary reports no located general affirmative answer through 2026 and makes an uncited non-biautomaticity assertion for \(g\ge2\), whose validity requires review.
\item Dataset attribution: Ulam AI / UnsolvedMath Contributors, CC BY 4.0. Source rights NEEDS\_REVIEW: publicly accessible author/publisher copy, redistribution terms requiring review.
\end{itemize}
\end{document}
